\documentclass[11pt]{article}
\usepackage{mathunal-base}
\mathunalfuente{serif}
\newcommand{\resp}{\par\vspace{2pt}\noindent\underline{\hspace{9cm}}\par\vspace{6pt}}

\begin{document}

{\large\bfseries Matemáticas Básicas --- Modelo de examen final}

\medskip
\noindent Escuela de Matemáticas. Universidad Nacional de Colombia, Sede Medellín.

\medskip
\noindent\textit{Duración: 1 hora y 50 minutos. Verifique que sus datos de identificación sean correctos. No desengrape su examen ni añada otras grapas. No está permitido: hacer preguntas, usar calculadoras ni trabajar en hojas diferentes a las del examen. No se permite que ningún celular se encuentre a la vista. Diligencie sus respuestas en la tabla correspondiente: escriba grande y claro. Cada pregunta vale 10\,\%.}

\begin{enumerate}
\item $[10\%]$ El resultado de factorizar completamente sobre $\mathbb{R}$ la expresión $25x^4 - 49y^{16}$ es:
\resp

\item $[10\%]$ El conjunto solución de la ecuación
\[
\sqrt{5 - x} - 1 = 2 - x
\]
es:
\resp

\item $[10\%]$ La ecuación de la recta que pasa por el punto $\left(-\tfrac{1}{3},\ \tfrac{11}{4}\right)$ y es perpendicular a la recta que pasa por los puntos $(-1, 3)$ y $(2, -1)$ es:
\begin{opciones}
\item $y = \tfrac{3}{4}x - 1$
\item $y = \tfrac{3}{4}x + 3$
\item $y = \tfrac{3}{4}x + 1$
\item $y = \tfrac{3}{4}x + 2$
\end{opciones}

\item $[10\%]$ Racionalizando el numerador de $\dfrac{\sqrt5 + \sqrt w}{w}$ y simplificando se obtiene:
\resp

\item $[10\%]$ Al realizar la operación
\[
\frac{x}{x^2 + x - 6} - \frac{2}{x^2 + 3x - 10}
\]
y simplificar, se obtiene:
\begin{opciones}
\item $\dfrac{x^2 + 3x - 6}{(x+3)(x-2)(x+5)}$
\item $\dfrac{x^2 + 2x - 9}{(x+3)(x-2)(x+5)}$
\item $\dfrac{x^2 + 2x - 5}{(x+3)(x-2)(x+5)}$
\item $\dfrac{x^2 - 2x + 1}{(x+3)(x-2)(x+5)}$
\end{opciones}

\item $[10\%]$ El residuo que se obtiene al dividir el polinomio $x^{34} + 7x^{13} - 2x^2 - x + 7$ entre el polinomio $x + 1$ es:
\resp

\item $[10\%]$ El conjunto de los ceros reales del polinomio $P(x) = x^3 + x^2 - 6x + 24$ es:
\resp

\item $[10\%]$ Los valores que las constantes $B$ y $C$ deben tener para que el par $(x, y) = (1, 2)$ sea solución del sistema de ecuaciones
\[
\begin{cases}
(C + 4)x + By = 11 \\
(B - 4)x + Cy = -5
\end{cases}
\]
son:
\resp

\item $[10\%]$ El centro de la circunferencia con ecuación $x^2 + y^2 - 10x - 8y + 32 = 0$ es:
\resp

\item $[10\%]$ La solución del sistema
\[
\begin{cases}
-9x + 2y = 29 \\
18x - 5y = -59
\end{cases}
\]
es:
\resp
\end{enumerate}

\vfill
\noindent{\small\itshape Transcripción digital --- MathUNAL}

\end{document}
